\section{Primary vertex}
\label{sec:ch6_primary_vertex}

\subsection{Reconstruction}

Collision vertices are reconstructed from tracks compatible with a common origin~\cite{ATLAS2015Vertex}.
Candidate positions are first identified from the longitudinal distribution of tracks near the beam line.
The vertex positions are then fitted using the track parameters and their uncertainties, with reduced weights assigned to incompatible tracks.
Tracks that are not assigned to a fitted vertex are considered when searching for further vertices.
Several collision vertices can therefore be reconstructed in an event containing pile-up.

\subsection{Primary-vertex selection}

The primary vertex associated with the hard scattering is selected from vertices with at least two tracks satisfying $\pT>500~\MeV$~\cite{ATLAS2026LQTauNu}.
The vertex with the largest sum of squared track transverse momenta is chosen:
\begin{align}
    v_{\mathrm{PV}}=\operatorname*{arg\,max}_{v}
    \sum_{i\in v}p_{\mathrm{T},i}^{\,2}.
    \label{eq:ch6_primary_vertex}
\end{align}
Energetic tracks are given greater weight by this criterion, favouring the hard scattering over additional soft collisions.
Events without a qualifying vertex are rejected.
The selected vertex is used as the reference for prompt-lepton selection, isolation, pile-up jet rejection and the track contribution to $\met$.
